% 第 14 章 Vulkan 后端迁移（占位——全书最重章，素材最富）
锚点：\texttt{Platform/Vulkan/} 40 文件约 6000 行（Context/Swapchain/
Allocator(VMA)/RenderPass/Pipeline/PipelineCache/Descriptor 三层/Shader
(shaderc+spirv-cross)/Buffer/Texture/VertexArray/Synchronization/
BarrierUtil/AsyncReadback/IBLGenerator）。

叙事骨架（SPEC.md Phase 1-8 即章节骨架）：
Phase 1 抽象泄漏归零（前置工程）→ Phase 2 基础设施清屏 →
Phase 3 RendererAPI 核心 → Phase 4 shaderc 运行时编译 →
Phase 5 VMA 资源族 → Phase 6 ImGui 多视口 → Phase 7 Compute 全家桶迁移
（IBL/草地/网格/粒子/SPH）→ Phase 8 PBR 主路径接通（进行中，
DrawIndexed 待实装）。

决策素材：SPEC §3 十六条 D-x 决策全部按``背景+结论+影响范围''格式现成；
ADR-0001（shader 单源双路径三方案对比）/ADR-0002（主帧 cmd 暴露策略）
标准 MADR 格式；Pitfalls 表 17 条即本章``踩坑录''小节初稿。

% TODO(M3): 与 OpenGL 抽象层的映射表；双路径 shader 宏策略；验证方法学
%（validation layer + 冒烟层）。
